\documentclass[11pt]{article}
% Standalone preamble for web-only supplementary notes.
\usepackage[T1]{fontenc}
\usepackage{amsmath,amssymb,amsthm}
\usepackage{hyperref}
\usepackage{cleveref}
\DeclareMathOperator{\Map}{Map}
\DeclareMathOperator{\Fun}{Fun}
\newcommand{\iso}{\xrightarrow{\sim}}
\newtheorem{lemma}{Supplementary Lemma}
\newtheorem{corollary}[lemma]{Supplementary Corollary}
\crefname{lemma}{Supplementary Lemma}{Supplementary Lemmas}
\Crefname{lemma}{Supplementary Lemma}{Supplementary Lemmas}
\crefname{corollary}{Supplementary Corollary}{Supplementary Corollaries}
\Crefname{corollary}{Supplementary Corollary}{Supplementary Corollaries}

\begin{document}

\section*{The compatibility hidden in the short proof}

The book's \href{functor-categories.html\#lem:Functor_Category_Preserves_Pullbacks}{pullback-preservation lemma}
proves that $\Fun(T,-)$ preserves pullbacks by testing with $\Map(S,-)$ and
then uncurrying. The corner equivalences are
\[
  u_Z\colon \Map(S,\Fun(T,Z))\iso\Map(S\times T,Z).
\]
To use them, we must compare the specified commutativity identifications as
well as the four arrows. This note expands that step and explains how the
resulting comparison transfers the pullback property.

We work with the absolute categories and the mapping-anima and functor-category
structures of Sections 1.1--1.7. The cone conversion below does not require a
pullback hypothesis. The later application also uses the pullback structure of
Section 1.6. Statements about parameters $X$ below explicitly require $X$ to be
an anima.

\section*{Keeping the commutativity data}

For a cospan $C\xrightarrow{f}E\xleftarrow{g}D$, a cone on $X$ consists of
$a\colon X\to C$, $b\colon X\to D$, and an identification
$\theta\colon f\circ a=g\circ b$. A
%!% begin supp-cone-comparison
comparison of cones
%!% end supp-cone-comparison
from $(a,b,\theta)$ to $(a',b',\theta')$ consists of identifications
$\alpha\colon a=a'$ and $\beta\colon b=b'$, together with an identification
\[
  \theta'\cdot(f\circ\alpha)=(g\circ\beta)\cdot\theta.
\]
Here a dot denotes composition of identifications, with the right-hand
identification applied first. Thus the comparison remembers a filling of the
square between the two matchings. Precomposition of a cone uses the specified
associativity identifications to give its new matching. We write $h^*q$ for
this cone and $q\simeq q'$ for a comparison of cones in this sense.

Fix categories $S,T$ and put
\[
 A_Z=\Map(S,\Fun(T,Z)),\qquad B_Z=\Map(S\times T,Z).
\]
Write $f_A,g_A$ and $f_B,g_B$ for the induced postcomposition functors on
these two families. These symbols name the cospans used below; we do not need
to construct a category of all cones.

\begin{lemma}
\label[lemma]{lem:Supplementary_Cone_Conversion}
For every anima $X$, there are conversions $U_X$ and $V_X$ between cones on
$X$ over $(f_A,g_A)$ and over $(f_B,g_B)$. For a cone $r=(a,b,\theta)$,
the legs of $U_X(r)$ are $u_C\circ a$ and $u_D\circ b$.
Both conversions take comparisons of cones to comparisons of cones, and there
are comparisons $V_X(U_X(r))\simeq r$ and $U_X(V_X(t))\simeq t$.
Moreover, a comparison $U_X(r)\simeq U_X(r')$ gives a comparison $r\simeq r'$.
\end{lemma}
\begin{proof}[Proof idea]
Uncurry the legs and matching of $r$ twice. This gives a cone over $(f,g)$ on
$(X\times S)\times T$. Restrict along the associativity equivalence
$X\times(S\times T)\to(X\times S)\times T$, and curry into the mapping
anima for $S\times T$. The resulting legs are identified with
$u_C\circ a$ and $u_D\circ b$. Transport the matching along these two
identifications to obtain exactly these legs.

For the reverse conversion, uncurry a cone over $(f_B,g_B)$, change the
parenthesization in the opposite direction, and curry first into the functor
category and then into the mapping anima. The evaluation comparisons and
the inverse comparisons for product associativity give the two round trips.
These are comparisons of whole cones, so they retain compatibility with the
matchings. Uncurrying comparisons and reflecting them through the relevant
mapping equivalences also gives the final reflection assertion.
\end{proof}

\section*{Comparing the specified squares}

Let $q=(p,r,\eta)$ be any cone over $(f,g)$ with vertex $W$. Denote the cones
obtained by the specified postcomposition constructions by $A(q)$ and $B(q)$,
with vertices $A_W$ and $B_W$ respectively.

\begin{lemma}
\label[lemma]{lem:Supplementary_Square_Compatibility}
For every anima $X$ and every $h\colon X\to A_W$, there is a comparison
\[
 U_X\bigl(h^*A(q)\bigr)\simeq(u_W\circ h)^*B(q).
\]
In particular, converting $A(q)$ gives a cone comparison with $u_W^*B(q)$.
This assertion does not require $q$ to be a pullback cone.
\end{lemma}
\begin{proof}[Proof idea]
Double uncurrying and reassociation turn $h$ into a functor
$\bar h\colon X\times(S\times T)\to W$. The evaluation comparison for
the specified mapped cone identifies the uncurried left-hand cone with
$\bar h^*q$. Uncurrying $u_W\circ h$ gives a second functor into $W$,
equipped with the comparison expressing that $u_W$ represents uncurrying.
Acting on $q$ by this comparison compares the two induced cones, including
their matching identifications. Reflection through the mapping equivalence
gives the asserted cone comparison.

More explicitly, let $\alpha$ and $\beta$ be its two leg comparisons and
let $\mu$ and $\nu$ be the source and target matchings. The compatibility
supplied by the comparison is
\[
 \nu\cdot(f_B\circ\alpha)=(g_B\circ\beta)\cdot\mu.
\]
Equivalently, after transporting $\mu$ along $\alpha$ and $\beta$, it is
identified with $\nu$. This is the precise matching obligation. Taking
$X=A_W$ and $h$ the identity, and inserting the unit comparisons, gives the
assertion about the square itself. The category $A_W$ is an anima.
\end{proof}

The leg comparisons here are those constructed by this cone conversion.
The argument does not require an additional identification with independently
chosen postcomposition-naturality witnesses. Nor does it assert all higher
coherence laws for a category of cone diagrams.

\section*{Transferring the pullback property}

\begin{lemma}
\label[lemma]{lem:Supplementary_Pullback_Transfer}
Let $s$ and $t$ be cones over two cospans, with vertices $Y$ and $Z$, and let
$u\colon Y\to Z$ be an equivalence. Suppose that, for every anima $X$, a
conversion $K_X$ between cones over the two cospans preserves and reflects
cone comparisons, and is equipped, for every $h\colon X\to Y$, with a comparison
\[
 K_X(h^*s)\simeq(u\circ h)^*t.
\]
If $t$ is a pullback cone, $Y$ is an anima, and the chosen pullback of the
source cospan is an anima, then $s$ is a pullback cone.
\end{lemma}
\begin{proof}[Proof idea]
For a cone $r$ on an anima $X$ over the source cospan, factor $K_X(r)$
through $t$. Lift the factorizing functor through $u$. The displayed
comparison and reflection of cone comparisons identify the cone induced by
this lift with $r$.

If two functors $h,k\colon X\to Y$ induce comparable cones, applying
$K_X$ and the displayed comparisons gives comparable cones induced by
$u\circ h$ and $u\circ k$. The pullback property of $t$ gives an
identification between these functors, which reflects through $u$.

Apply the first construction to the chosen pullback cone of the source
cospan. It supplies a candidate inverse to the comparison from $Y$ to that
pullback. Its first inverse comparison follows from its cone computation
and the universal property of the chosen pullback. The second follows from
the reflection property just proved, applied with $X=Y$. These two tests
are legitimate by the anima hypotheses and prove that the comparison is an
equivalence.
\end{proof}

\begin{corollary}
\label[corollary]{cor:Supplementary_Functor_Pullbacks}
For every category $T$, postcomposition sends a pullback cone $q$ to a
pullback cone under $\Fun(T,-)$, with its specified commutativity identification.
\end{corollary}
\begin{proof}[Proof idea]
For every category $S$, the cone $B(q)$ is a pullback by the mapping-anima
criterion in \href{pullbacks.html}{Section 1.6}. Apply
\Cref{lem:Supplementary_Pullback_Transfer} using $u_W$, the conversion of
\Cref{lem:Supplementary_Cone_Conversion}, and the comparison of
\Cref{lem:Supplementary_Square_Compatibility}. Its anima hypotheses hold:
$A_W$ is a mapping anima, and pullbacks of animae are animae. Thus $A(q)$
is a pullback for every $S$. Mapping-anima detection now shows that the
specified $\Fun(T,-)$ cone is a pullback. This recovers the short argument
in the book with its compatibility step made explicit.
\end{proof}

\end{document}
